% Paper D: Tabling and Mode-Driven Goal Scheduling
% The "Termination Paper" - SLG tabling + mode analysis = termination + optimal scheduling
%
% For God so loved the world that he gave his only begotten Son,
% that whoever believes in him should not perish but have eternal life.
% John 3:16

\documentclass[11pt,a4paper]{article}

\usepackage[utf8]{inputenc}
\usepackage{amsmath,amssymb,amsthm}
\usepackage{booktabs}
\usepackage{hyperref}
\usepackage{xspace}
\usepackage{listings}
\usepackage{algorithm}
\usepackage{algpseudocode}
\usepackage{tikz}
\usetikzlibrary{arrows.meta,positioning,shapes}

% Theorem environments
\newtheorem{theorem}{Theorem}
\newtheorem{lemma}[theorem]{Lemma}
\newtheorem{definition}[theorem]{Definition}
\newtheorem{proposition}[theorem]{Proposition}
\newtheorem{corollary}[theorem]{Corollary}

% Include shared definitions
% Shared Definitions for miniKanren 1-Bit Paper Suite
% Include with: % Shared Definitions for miniKanren 1-Bit Paper Suite
% Include with: % Shared Definitions for miniKanren 1-Bit Paper Suite
% Include with: \input{../shared_chirho/shared_defs_chirho}
%
% For God so loved the world that he gave his only begotten Son,
% that whoever believes in him should not perish but have eternal life.
% John 3:16

% ============================================================================
% Core Definitions (shared across all papers)
% ============================================================================

% The Finite-Domain miniKanren Kernel
\newcommand{\fdmk}{FD-miniKanren\xspace}  % Finite-Domain miniKanren

% ============================================================================
% Scope Clarification (include in each paper's introduction)
% ============================================================================

\newcommand{\scopeclarification}{%
\textbf{Scope Clarification.}
This paper addresses the \emph{Finite-Domain miniKanren Kernel} (\fdmk),
where variable domains are finite bitmasks.
Standard miniKanren unification over unbounded term structures uses a separate
reference implementation; the bit-parallel approach accelerates the constraint
propagation core, not full structural unification.
}

% ============================================================================
% Common Notation
% ============================================================================

% Domains
\newcommand{\dom}[1]{D_{#1}}          % Domain of variable x
\newcommand{\domand}{\land}            % Domain intersection (AND)
\newcommand{\domor}{\lor}              % Domain union (OR)

% Search
\newcommand{\state}{(\vec{D}, \sigma)} % Search state
\newcommand{\findChirho}{\text{find}}      % Union-find find
\newcommand{\unionChirho}{\text{union}}    % Union-find union

% Tensors
\newcommand{\tensor}[1]{T_{#1}}        % Tensor for relation R
\newcommand{\contract}{\circledast}    % Tensor contraction

% Semirings
\newcommand{\softand}{\odot}           % Soft AND (multiply)
\newcommand{\softor}{\oplus}           % Soft OR (prob sum)

% ============================================================================
% Hierarchical Domain Scaling
% ============================================================================

% Domain type naming (semantic macros for consistency across papers)
% Using semantic names to avoid confusion with Rust type identifiers
\newcommand{\domainSmallChirho}{\texttt{BitVec64Chirho}}
\newcommand{\domainMediumChirho}{\texttt{Hierarchical4kChirho}}
\newcommand{\domainLargeChirho}{\texttt{Hierarchical256kChirho}}
\newcommand{\domainDiffChirho}{\texttt{DiffHierarchical4kChirho}}

% Measured performance values (from Rust benchmarks, Jan 2026)
\newcommand{\domainSmallTimeChirho}{420 ps}
\newcommand{\domainMediumTimeChirho}{39 ns}
\newcommand{\domainLargeTimeChirho}{2.5 $\mu$s}
\newcommand{\domainLargeOpsChirho}{400K ops/sec}

% ============================================================================
% Paper Suite Reference
% ============================================================================

% Use these to cite sibling papers
\newcommand{\paperA}{Paper A: Bit-Parallel Finite-Domain Relational Search}
\newcommand{\paperB}{Paper B: Relations as Sparse Boolean Tensors}
\newcommand{\paperC}{Paper C: Bridging Infinite Domains with Hash-Consed Term IDs}
\newcommand{\paperD}{Paper D: Tabling and Mode-Driven Goal Scheduling}
\newcommand{\paperE}{Paper E: Fully Fused Logic Accelerator}
\newcommand{\paperF}{Paper F: Differentiable Constraint Solving via Semiring Relaxation}

% ============================================================================
% Acknowledgments Template
% ============================================================================

\newcommand{\standardacknowledgments}{%
Above all, I thank my Lord and Savior Jesus Christ for saving a wretch like me.

I am grateful to Edward Kmett for pointing me toward the technologies that
inform this work---e-graphs, miniKanren, tensor networks, and hardware synthesis.

I also thank my father, Roger Cowan, for his patient support.

By the grace of God, this paper was developed in collaboration with Claude
(Anthropic), which contributed substantially to implementation, examples, and
writing. Google Gemini and OpenAI's GPT-5.2 Codex provided valuable feedback
on drafts.
}

% ============================================================================
% Example environment (for papers that use it)
% ============================================================================
\newenvironment{exampleChirho}{\par\noindent\textbf{Example.}\itshape}{\par}

% Soli Deo Gloria

%
% For God so loved the world that he gave his only begotten Son,
% that whoever believes in him should not perish but have eternal life.
% John 3:16

% ============================================================================
% Core Definitions (shared across all papers)
% ============================================================================

% The Finite-Domain miniKanren Kernel
\newcommand{\fdmk}{FD-miniKanren\xspace}  % Finite-Domain miniKanren

% ============================================================================
% Scope Clarification (include in each paper's introduction)
% ============================================================================

\newcommand{\scopeclarification}{%
\textbf{Scope Clarification.}
This paper addresses the \emph{Finite-Domain miniKanren Kernel} (\fdmk),
where variable domains are finite bitmasks.
Standard miniKanren unification over unbounded term structures uses a separate
reference implementation; the bit-parallel approach accelerates the constraint
propagation core, not full structural unification.
}

% ============================================================================
% Common Notation
% ============================================================================

% Domains
\newcommand{\dom}[1]{D_{#1}}          % Domain of variable x
\newcommand{\domand}{\land}            % Domain intersection (AND)
\newcommand{\domor}{\lor}              % Domain union (OR)

% Search
\newcommand{\state}{(\vec{D}, \sigma)} % Search state
\newcommand{\findChirho}{\text{find}}      % Union-find find
\newcommand{\unionChirho}{\text{union}}    % Union-find union

% Tensors
\newcommand{\tensor}[1]{T_{#1}}        % Tensor for relation R
\newcommand{\contract}{\circledast}    % Tensor contraction

% Semirings
\newcommand{\softand}{\odot}           % Soft AND (multiply)
\newcommand{\softor}{\oplus}           % Soft OR (prob sum)

% ============================================================================
% Hierarchical Domain Scaling
% ============================================================================

% Domain type naming (semantic macros for consistency across papers)
% Using semantic names to avoid confusion with Rust type identifiers
\newcommand{\domainSmallChirho}{\texttt{BitVec64Chirho}}
\newcommand{\domainMediumChirho}{\texttt{Hierarchical4kChirho}}
\newcommand{\domainLargeChirho}{\texttt{Hierarchical256kChirho}}
\newcommand{\domainDiffChirho}{\texttt{DiffHierarchical4kChirho}}

% Measured performance values (from Rust benchmarks, Jan 2026)
\newcommand{\domainSmallTimeChirho}{420 ps}
\newcommand{\domainMediumTimeChirho}{39 ns}
\newcommand{\domainLargeTimeChirho}{2.5 $\mu$s}
\newcommand{\domainLargeOpsChirho}{400K ops/sec}

% ============================================================================
% Paper Suite Reference
% ============================================================================

% Use these to cite sibling papers
\newcommand{\paperA}{Paper A: Bit-Parallel Finite-Domain Relational Search}
\newcommand{\paperB}{Paper B: Relations as Sparse Boolean Tensors}
\newcommand{\paperC}{Paper C: Bridging Infinite Domains with Hash-Consed Term IDs}
\newcommand{\paperD}{Paper D: Tabling and Mode-Driven Goal Scheduling}
\newcommand{\paperE}{Paper E: Fully Fused Logic Accelerator}
\newcommand{\paperF}{Paper F: Differentiable Constraint Solving via Semiring Relaxation}

% ============================================================================
% Acknowledgments Template
% ============================================================================

\newcommand{\standardacknowledgments}{%
Above all, I thank my Lord and Savior Jesus Christ for saving a wretch like me.

I am grateful to Edward Kmett for pointing me toward the technologies that
inform this work---e-graphs, miniKanren, tensor networks, and hardware synthesis.

I also thank my father, Roger Cowan, for his patient support.

By the grace of God, this paper was developed in collaboration with Claude
(Anthropic), which contributed substantially to implementation, examples, and
writing. Google Gemini and OpenAI's GPT-5.2 Codex provided valuable feedback
on drafts.
}

% ============================================================================
% Example environment (for papers that use it)
% ============================================================================
\newenvironment{exampleChirho}{\par\noindent\textbf{Example.}\itshape}{\par}

% Soli Deo Gloria

%
% For God so loved the world that he gave his only begotten Son,
% that whoever believes in him should not perish but have eternal life.
% John 3:16

% ============================================================================
% Core Definitions (shared across all papers)
% ============================================================================

% The Finite-Domain miniKanren Kernel
\newcommand{\fdmk}{FD-miniKanren\xspace}  % Finite-Domain miniKanren

% ============================================================================
% Scope Clarification (include in each paper's introduction)
% ============================================================================

\newcommand{\scopeclarification}{%
\textbf{Scope Clarification.}
This paper addresses the \emph{Finite-Domain miniKanren Kernel} (\fdmk),
where variable domains are finite bitmasks.
Standard miniKanren unification over unbounded term structures uses a separate
reference implementation; the bit-parallel approach accelerates the constraint
propagation core, not full structural unification.
}

% ============================================================================
% Common Notation
% ============================================================================

% Domains
\newcommand{\dom}[1]{D_{#1}}          % Domain of variable x
\newcommand{\domand}{\land}            % Domain intersection (AND)
\newcommand{\domor}{\lor}              % Domain union (OR)

% Search
\newcommand{\state}{(\vec{D}, \sigma)} % Search state
\newcommand{\findChirho}{\text{find}}      % Union-find find
\newcommand{\unionChirho}{\text{union}}    % Union-find union

% Tensors
\newcommand{\tensor}[1]{T_{#1}}        % Tensor for relation R
\newcommand{\contract}{\circledast}    % Tensor contraction

% Semirings
\newcommand{\softand}{\odot}           % Soft AND (multiply)
\newcommand{\softor}{\oplus}           % Soft OR (prob sum)

% ============================================================================
% Hierarchical Domain Scaling
% ============================================================================

% Domain type naming (semantic macros for consistency across papers)
% Using semantic names to avoid confusion with Rust type identifiers
\newcommand{\domainSmallChirho}{\texttt{BitVec64Chirho}}
\newcommand{\domainMediumChirho}{\texttt{Hierarchical4kChirho}}
\newcommand{\domainLargeChirho}{\texttt{Hierarchical256kChirho}}
\newcommand{\domainDiffChirho}{\texttt{DiffHierarchical4kChirho}}

% Measured performance values (from Rust benchmarks, Jan 2026)
\newcommand{\domainSmallTimeChirho}{420 ps}
\newcommand{\domainMediumTimeChirho}{39 ns}
\newcommand{\domainLargeTimeChirho}{2.5 $\mu$s}
\newcommand{\domainLargeOpsChirho}{400K ops/sec}

% ============================================================================
% Paper Suite Reference
% ============================================================================

% Use these to cite sibling papers
\newcommand{\paperA}{Paper A: Bit-Parallel Finite-Domain Relational Search}
\newcommand{\paperB}{Paper B: Relations as Sparse Boolean Tensors}
\newcommand{\paperC}{Paper C: Bridging Infinite Domains with Hash-Consed Term IDs}
\newcommand{\paperD}{Paper D: Tabling and Mode-Driven Goal Scheduling}
\newcommand{\paperE}{Paper E: Fully Fused Logic Accelerator}
\newcommand{\paperF}{Paper F: Differentiable Constraint Solving via Semiring Relaxation}

% ============================================================================
% Acknowledgments Template
% ============================================================================

\newcommand{\standardacknowledgments}{%
Above all, I thank my Lord and Savior Jesus Christ for saving a wretch like me.

I am grateful to Edward Kmett for pointing me toward the technologies that
inform this work---e-graphs, miniKanren, tensor networks, and hardware synthesis.

I also thank my father, Roger Cowan, for his patient support.

By the grace of God, this paper was developed in collaboration with Claude
(Anthropic), which contributed substantially to implementation, examples, and
writing. Google Gemini and OpenAI's GPT-5.2 Codex provided valuable feedback
on drafts.
}

% ============================================================================
% Example environment (for papers that use it)
% ============================================================================
\newenvironment{exampleChirho}{\par\noindent\textbf{Example.}\itshape}{\par}

% Soli Deo Gloria


% Code listings style
\lstset{
  basicstyle=\ttfamily\small,
  keywordstyle=\bfseries,
  commentstyle=\itshape,
  columns=flexible,
  breaklines=true,
  frame=single,
  numbers=left,
  numberstyle=\tiny
}

\title{Tabling and Mode-Driven Goal Scheduling\\
\large The Termination Paper}

\author{
  L.J. Brian Cowan \\
  \texttt{loveJesus@loveJes.us}
}

\date{Draft - \today}

\begin{document}

\maketitle

% ============================================================================
\begin{abstract}
% ============================================================================

We present a tabling system that guarantees termination for finite-domain
relational queries through mode-driven goal scheduling. The key insight is that
\emph{goal reordering by groundness} transforms potentially non-terminating
left-to-right evaluation into terminating schedules.

Given a conjunction $(g_1, g_2, \ldots, g_n)$, we dynamically reorder goals
so that those with more ground arguments execute first. This ensures that
generators (which enumerate infinite streams) are deferred until constraining
goals have bound their shared variables.

For mutually recursive predicates, we employ SLG resolution with strongly
connected component (SCC) detection. Goals within an SCC must complete
together; Tarjan's algorithm identifies these groups in $O(V+E)$ time.

Our implementation demonstrates:
\begin{itemize}
  \item \textbf{Termination}: Composition queries that would loop without
        tabling now terminate with finite answer sets
  \item \textbf{Optimal scheduling}: Mode analysis achieves automatic goal
        reordering equivalent to hand-tuned Prolog
  \item \textbf{Tensor construction}: Each solved query populates a sparse
        Boolean tensor incrementally, bridging streams and matrices
\end{itemize}

We prove that for \fdmk with bounded term depth, tabling with mode analysis
guarantees termination while preserving completeness.

\end{abstract}

% ============================================================================
\section{Introduction}
\label{sec:intro_chirho}
% ============================================================================

\scopeclarification

Consider the query:
\begin{verbatim}
  (run* (A B X C)
    (appendo A B X)
    (appendo X C [0 1 2]))
\end{verbatim}

Evaluated left-to-right, this query diverges. The first goal
\texttt{(appendo A B X)} generates an infinite stream of values for $X$,
never reaching the second goal that would constrain it.

Yet the query has exactly 10 solutions. How can we find them?

The answer lies in \emph{mode-driven goal reordering}. By analyzing which
arguments are ground (known) versus free (unknown), we can reorder goals
to ensure termination:

\begin{enumerate}
  \item \texttt{(appendo X C [0 1 2])} --- output \texttt{[0 1 2]} is ground,
        so this runs in \emph{backward mode}, enumerating finite splits
  \item \texttt{(appendo A B X)} --- now $X$ is bound, so this also runs
        in backward mode
\end{enumerate}

This paper formalizes this transformation and proves its correctness.

\paragraph{Contributions.}
\begin{enumerate}
  \item \textbf{Goal reordering by groundness} (\S\ref{sec:reordering_chirho}):
        A dynamic scheduling algorithm that reorders conjunctions based on
        argument groundness
  \item \textbf{Tabling as incremental tensor construction}
        (\S\ref{sec:tensor_chirho}): Each tabled answer contributes a 1-bit
        entry to a sparse Boolean tensor
  \item \textbf{SCC detection for mutual recursion} (\S\ref{sec:scc_chirho}):
        Tarjan's algorithm identifies mutually recursive goal clusters that
        must complete together
\end{enumerate}

% ============================================================================
\section{Background: miniKanren and Tabling}
\label{sec:background_chirho}
% ============================================================================

\subsection{miniKanren Semantics}
\label{subsec:mk_chirho}

miniKanren is a relational programming language where programs are
\emph{relations} rather than functions. The core operators are:

\begin{itemize}
  \item \texttt{(== t1 t2)} --- unification goal
  \item \texttt{(conde [g ...] ...)} --- disjunction with interleaving
  \item \texttt{(fresh (x ...) g ...)} --- introduce fresh variables
\end{itemize}

Evaluation produces a \emph{stream} of substitutions, potentially infinite.

\subsection{The Termination Problem}
\label{subsec:termination_chirho}

Without tabling, recursive relations can diverge:

\begin{lstlisting}[caption={The appendo relation}]
(defrel (appendo l s out)
  (conde
    [(== l '()) (== s out)]
    [(fresh (h t res)
       (== l (cons h t))
       (== out (cons h res))
       (appendo t s res))]))
\end{lstlisting}

Query \texttt{(run* (q) (appendo q '(1 2) '(0 1 2)))} terminates.

Query \texttt{(run* (q) (appendo q q q))} diverges.

The difference is \emph{mode}: the first has a ground output (backward mode),
while the second has no ground arguments (generate mode).

\subsection{SLG Resolution}
\label{subsec:slg_chirho}

SLG (Selective Linear Definite clause resolution with Ground negation) is a
tabling strategy that:

\begin{enumerate}
  \item Memoizes subgoal answers in a table
  \item Suspends when encountering a variant call
  \item Resumes suspended computations when new answers arrive
  \item \emph{Completes} goals when no more answers are possible
\end{enumerate}

Completion is the key innovation: a goal is complete when all its answers
have been found. For mutually recursive goals, completion must occur
simultaneously for all goals in the same strongly connected component.

% ============================================================================
\section{Contribution 1: Goal Reordering by Groundness}
\label{sec:reordering_chirho}
% ============================================================================

\subsection{Groundness Scoring}
\label{subsec:scoring_chirho}

For a goal $g$ with arguments $(a_1, \ldots, a_k)$ under substitution $\sigma$,
define the \emph{groundness score}:

\begin{definition}[Groundness Score]
\label{def:groundness_chirho}
\[
\text{score}_\sigma(g) = \sum_{i=1}^{k} w(a_i, \sigma)
\]
where
\[
w(a, \sigma) = \begin{cases}
  10 & \text{if } \text{walk}^*(a, \sigma) \text{ is ground} \\
  5  & \text{if } \text{walk}^*(a, \sigma) \text{ is partially ground} \\
  0  & \text{otherwise}
\end{cases}
\]
\end{definition}

\subsection{Dynamic Reordering Algorithm}
\label{subsec:algorithm_chirho}

\begin{algorithm}
\caption{Smart Conjunction (\texttt{smart\_conj\_chirho})}
\label{alg:smart_conj_chirho}
\begin{algorithmic}[1]
\Procedure{SmartConj}{$goals$, $state$}
  \If{$|goals| = 0$}
    \State \textbf{yield} $state$
  \ElsIf{$|goals| = 1$}
    \State \textbf{yield from} $goals[0](state)$
  \Else
    \State $sorted \gets \text{sort}(goals, \text{key}=\lambda g. -\text{score}_{state.\sigma}(g))$
    \State $first \gets sorted[0]$
    \State $rest \gets sorted[1:]$
    \For{$st' \in first(state)$}
      \State \textbf{yield from} \Call{SmartConj}{$rest$, $st'$}
    \EndFor
  \EndIf
\EndProcedure
\end{algorithmic}
\end{algorithm}

\subsection{Termination Guarantee}
\label{subsec:termination_guarantee_chirho}

\begin{theorem}[Reordering Termination]
\label{thm:reordering_termination_chirho}
For \fdmk with bounded term depth $d$, if a conjunction contains at least
one goal with a ground argument, smart reordering guarantees that all goals
eventually execute in backward or forward mode.
\end{theorem}

\begin{proof}
By induction on the number of free variables. Each backward-mode goal binds
at least one variable, reducing the set of free variables. Since the domain
is finite and term depth is bounded, this process terminates.
\end{proof}

% ============================================================================
\section{Contribution 2: Tabling as Incremental Tensor Construction}
\label{sec:tensor_chirho}
% ============================================================================

\subsection{The Tensor Store}
\label{subsec:tensor_store_chirho}

Each relation corresponds to a sparse Boolean tensor. For \texttt{appendo}:

\begin{definition}[Appendo Tensor]
\label{def:appendo_tensor_chirho}
\[
T_{\text{appendo}} \in \{0,1\}^{|L| \times |S| \times |O|}
\]
where $L$ is the set of possible first arguments, $S$ the second, and $O$
the output. Entry $T[l, s, o] = 1$ iff $l \mathbin{++} s = o$.
\end{definition}

\subsection{Incremental Population}
\label{subsec:incremental_chirho}

During query evaluation, each ground answer $(l, s, o)$ adds an entry:

\begin{lstlisting}[caption={Tensor population from tabling}]
def add_appendo_chirho(l, s, out):
    """Add a ground appendo triple to tensor"""
    l_list = term_to_pylist_chirho(l)
    s_list = term_to_pylist_chirho(s)
    out_list = term_to_pylist_chirho(out)

    if all_ground:
        TENSOR_STORE_CHIRHO.appendo_triples_chirho.add(
            (tuple(l_list), tuple(s_list), tuple(out_list))
        )
\end{lstlisting}

\subsection{The Bridge}
\label{subsec:bridge_chirho}

\begin{figure}[h]
\centering
\begin{tikzpicture}[
  box/.style={rectangle, draw, minimum width=3cm, minimum height=1cm, align=center},
  arrow/.style={-{Stealth[scale=1.2]}, thick}
]
  \node[box] (stream) at (0,0) {Lazy Streams\\(miniKanren)};
  \node[box] (table) at (5,0) {Tabling\\(bridge)};
  \node[box] (tensor) at (10,0) {Sparse Tensors\\(1-bit matrix)};

  \draw[arrow] (stream) -- (table) node[midway,above] {SLG};
  \draw[arrow] (table) -- (tensor) node[midway,above] {COO};
\end{tikzpicture}
\caption{Tabling bridges lazy streams and eager tensors}
\label{fig:bridge_chirho}
\end{figure}

% ============================================================================
\section{Contribution 3: SCC Detection for Mutual Recursion}
\label{sec:scc_chirho}
% ============================================================================

\subsection{The Mutual Recursion Problem}
\label{subsec:mutual_problem_chirho}

Consider the classic even/odd predicates:

\begin{lstlisting}[caption={Mutually recursive even/odd}]
even(0).
even(s(X)) :- odd(X).
odd(s(X)) :- even(X).
\end{lstlisting}

These predicates are \emph{mutually recursive}: \texttt{even} calls \texttt{odd}
and vice versa. Without proper handling, tabling can deadlock or miss answers.

\subsection{Dependency Graph and SCCs}
\label{subsec:dependency_chirho}

\begin{definition}[Goal Dependency Graph]
\label{def:dependency_graph_chirho}
The \emph{goal dependency graph} $G = (V, E)$ has:
\begin{itemize}
  \item $V$ = set of tabled goal patterns
  \item $(g_1, g_2) \in E$ iff evaluating $g_1$ may call $g_2$
\end{itemize}
\end{definition}

\begin{definition}[Strongly Connected Component]
\label{def:scc_chirho}
A \emph{strongly connected component} (SCC) is a maximal set of goals where
every goal is reachable from every other goal via the dependency edges.
\end{definition}

For even/odd, both predicates form a single SCC.

\subsection{Tarjan's Algorithm}
\label{subsec:tarjan_chirho}

We use Tarjan's algorithm to compute SCCs in $O(V+E)$ time:

\begin{lstlisting}[caption={SCC computation}]
def compute_sccs_chirho(self):
    """Compute SCCs using Tarjan's algorithm"""
    for pattern in self.goals_chirho:
        if pattern not in self.index_chirho:
            self._tarjan_chirho(pattern)

    # Assign SCC IDs
    for i, scc in enumerate(self.sccs_chirho):
        for pattern in scc:
            self.goals_chirho[pattern].scc_id_chirho = i
\end{lstlisting}

\subsection{Completion Protocol}
\label{subsec:completion_chirho}

\begin{theorem}[SCC Completion]
\label{thm:scc_completion_chirho}
For a set of mutually recursive goals forming an SCC, completion occurs
if and only if:
\begin{enumerate}
  \item No goal in the SCC has pending work
  \item All dependencies outside the SCC are complete
\end{enumerate}
\end{theorem}

This ensures that mutually recursive predicates produce complete answer sets.

% ============================================================================
\section{Implementation}
\label{sec:implementation_chirho}
% ============================================================================

\subsection{Artifacts}
\label{subsec:artifacts_chirho}

The implementation consists of two Python modules:

\begin{enumerate}
  \item \texttt{tabling\_complete\_chirho.py}: Mode-driven goal reordering with
        tabling for \texttt{appendo}
  \item \texttt{mutual\_recursion\_chirho.py}: SLG completion with SCC detection
        for even/odd
\end{enumerate}

Both are available at:
\url{https://github.com/loveJesus/minikanren_1bit_chirho}

\subsection{Key Data Structures}
\label{subsec:datastructures_chirho}

\begin{lstlisting}[caption={Goal information for scheduling}]
class GoalInfoChirho:
    """Information about a goal for scheduling"""
    def __init__(self, goal, args, name=""):
        self.goal_chirho = goal
        self.args_chirho = args
        self.name_chirho = name

    def groundness_score_chirho(self, subst):
        score = 0
        for arg in self.args_chirho:
            walked = walk_deep_chirho(arg, subst)
            if is_ground_chirho(walked):
                score += 10
            elif isinstance(walked, ConsChirho):
                score += 5  # Partially ground
        return score
\end{lstlisting}

% ============================================================================
\section{Evaluation}
\label{sec:evaluation_chirho}
% ============================================================================

\subsection{Termination Tests}
\label{subsec:termination_tests_chirho}

\begin{table}[h]
\centering
\begin{tabular}{lccc}
\toprule
Query & Without Reordering & With Reordering & Solutions \\
\midrule
\texttt{appendo(A,B,[0,1,2])} & Terminates & Terminates & 4 \\
\texttt{appendo(A,B,X), appendo(X,C,[0,1,2])} & Diverges & Terminates & 10 \\
Triple composition & Diverges & Terminates & 35 \\
\bottomrule
\end{tabular}
\caption{Termination comparison with and without goal reordering}
\label{tab:termination_chirho}
\end{table}

\subsection{Tensor Population}
\label{subsec:tensor_pop_chirho}

After running the composition query, the tensor store contains 15 ground
triples, forming a sparse representation of the appendo relation restricted
to the relevant domain.

\subsection{SCC Analysis}
\label{subsec:scc_analysis_chirho}

For the even/odd example with depth limit 15:
\begin{itemize}
  \item Goals tracked: 32
  \item SCCs detected: 16 (trivial) + 1 (even/odd mutual)
  \item Even tensor: \texttt{10101010...} (0, 2, 4, 6, ...)
  \item Odd tensor: \texttt{01010101...} (1, 3, 5, 7, ...)
\end{itemize}

% ============================================================================
\section{Related Work}
\label{sec:related_chirho}
% ============================================================================

\paragraph{Tabling in Prolog.}
XSB Prolog~\cite{swift1999} pioneered SLG resolution. Our work adapts these
ideas to the miniKanren setting with explicit mode analysis.

\paragraph{Mode Analysis.}
Mercury~\cite{somogyi1996} uses static mode declarations. We perform dynamic
mode inference at runtime based on substitution groundness.

\paragraph{Tensor Networks.}
The connection between logic programming and tensor networks has been explored
in Scallop~\cite{scallop2023}. We extend this to tabling and completion.

% ============================================================================
\section{Conclusion}
\label{sec:conclusion_chirho}
% ============================================================================

We have shown that tabling with mode-driven goal scheduling solves the
termination problem for finite-domain relational queries. The key insights:

\begin{enumerate}
  \item \textbf{Reordering = scheduling}: Dynamic goal reordering based on
        groundness is equivalent to optimal query planning
  \item \textbf{Tabling = tensor construction}: Each answer contributes a
        1-bit entry to a sparse Boolean tensor
  \item \textbf{SCC = fixpoint unit}: Mutually recursive goals must complete
        together, forming natural hardware compute units
\end{enumerate}

\paragraph{Code Availability.}
\url{https://github.com/loveJesus/minikanren_1bit_chirho}

\paragraph{Companion Papers.}
This paper is part of a suite:
\begin{itemize}
  \item \paperA
  \item \paperB
  \item \paperC
  \item \paperE
  \item \paperF
\end{itemize}

% ============================================================================
\section*{Acknowledgments}
% ============================================================================

\standardacknowledgments

\bibliographystyle{plain}
\begin{thebibliography}{9}

\bibitem{swift1999}
Swift, T., \& Warren, D. S. (1999).
\newblock An abstract machine for SLG resolution: Definite programs.
\newblock In {\em Proceedings of SLP'99}.

\bibitem{somogyi1996}
Somogyi, Z., Henderson, F., \& Conway, T. (1996).
\newblock The execution algorithm of Mercury: An efficient purely declarative
logic programming language.
\newblock {\em Journal of Logic Programming}, 29(1-3), 17--64.

\bibitem{scallop2023}
Li, Z., et al. (2023).
\newblock Scallop: From probabilistic Datalog to differentiable reasoning.
\newblock {\em NeurIPS 2023}.

\end{thebibliography}

\vspace{1em}
\begin{center}
\textit{Soli Deo Gloria} $\chi\rho$
\end{center}

\end{document}
